\section{REST e Panache como adaptadores}
\label{sec:panache-rest}

\subsection*{O resource REST é um adaptador de entrada}

Em \texttt{adaptador.entrada.rest}, o resource converte HTTP em chamada
de caso de uso e o resultado em DTO. Ele não tem regra de negócio:

\begin{lstlisting}[language=Java]
package br.com.turma.pedidos.adaptador.entrada.rest;
@Path("/pedidos")
public class PedidoResource {
    private final FecharPedido fecharPedido;
    @Inject
    PedidoResource(FecharPedido fecharPedido) { this.fecharPedido = fecharPedido; }

    @POST @Path("/{id}/fechamento")
    public PedidoDto fechar(@PathParam("id") String id) {
        return PedidoDto.de(fecharPedido.executar(new PedidoId(id)));
    }
}
\end{lstlisting}

\paragraph{Em que thread isso roda.}
O guia \citetitle{doc:quarkus-rest} diz que a thread depende da
assinatura do método. Métodos que devolvem \texttt{Uni},
\texttt{Multi}, \texttt{CompletionStage} ou \texttt{Publisher} são
considerados não bloqueantes e rodam na thread de I/O; os demais rodam
numa \emph{worker thread}. \texttt{@Blocking} e \texttt{@NonBlocking}
mudam o padrão, e um método com \texttt{@Transactional} é tratado como
bloqueante. O resource acima devolve um DTO comum, então roda numa worker
thread, e o caso de uso pode bloquear em JDBC sem travar o servidor. O
capítulo~\ref{cap:decisoes-quarkus} compara esse modelo com virtual
threads e com o reativo.

\subsection*{Panache: active record ou repository}

O guia \citetitle{doc:quarkus-panache} oferece dois estilos:
\begin{description}
\item[Active record] A entidade estende \texttt{PanacheEntity} e as
  consultas são métodos estáticos nela (\texttt{Pedido.findById(id)});
  o guia chama isso de \enquote{the Panache Active Record way}.
\item[Repository] Uma classe \texttt{@ApplicationScoped} implementa
  \texttt{PanacheRepository<PedidoEntity>} e oferece os mesmos métodos;
  estender \texttt{PanacheEntity} fica opcional.
\end{description}

O guia apresenta os dois sem recomendar um. Do ponto de vista da
Dependency Rule, a diferença é grande:

\begin{itemize}
\item \textbf{Active record coloca persistência dentro da entidade.} Se
  essa entidade for o seu \texttt{Pedido} de \texttt{dominio}, o círculo
  de dentro passa a importar \texttt{io.quarkus.hibernate.orm.panache} e
  \texttt{jakarta.persistence}, e uma regra de negócio pode chamar
  \texttt{persist()} no meio de um cálculo. A regra quebra no lugar mais
  central.
\item \textbf{Repository cabe num adaptador.} \texttt{PedidoEntity} (JPA)
  e \texttt{PedidoRepository} (Panache) ficam em
  \texttt{adaptador.saida.persistencia}; \texttt{PedidosPanache}
  implementa a porta \texttt{Pedidos} e converte entre
  \texttt{PedidoEntity} e o record \texttt{Pedido}.
\end{itemize}

\paragraph{O custo honesto.}
O mapeamento entidade$\leftrightarrow$domínio é código repetitivo, e
para um serviço que é essencialmente CRUD ele não protege nada: não há
regra de negócio para isolar. Nesse caso, active record é a escolha mais
simples e não é \enquote{errado}; é uma decisão consciente de que o
banco \emph{é} o modelo. Quando aparecer a primeira regra que não cabe
num \texttt{UPDATE}, extraia o domínio. Registre a escolha num ADR para
que o próximo dev saiba que foi intencional.

\paragraph{Testabilidade, de novo.}
O guia de Panache aponta mais uma diferença: com active record, o
Mockito não consegue simular os métodos estáticos, e é preciso usar o
\texttt{PanacheMock}; com repository, \enquote{you can use Mockito
directly}. Com a porta \texttt{Pedidos} no meio, nem isso: o teste do
caso de uso usa um fake em memória, como na Parte~\ref{parte:testes}.

\paragraph{O que o Panache faz no seu bytecode.}
Os getters e setters das entidades com campos \texttt{public} são
gerados no build (seção~\ref{sec:gizmo}). O guia descreve esse
\emph{enhancement} para as entidades; os records do domínio não são
entidades JPA, então continuam exatamente como o \texttt{javac} os
gerou (seção~\ref{sec:records}).
